\section{Conclusion}\label{sec:conclusion}

Six defensible occupational AI-exposure architectures construct different empirical treatments. They encode different technologies and occupational primitives, diverge in observable content and rankings, and draw identifying information from different occupations. Occupational mapping changes who enters the estimand, while common support does not eliminate differences in treatment-group membership or residual variation.

Despite that upstream divergence, every architecture produces a negative high-versus-low early-career employment-stock point estimate on literal common support. The confirmatory primary estimate also remains negative under every fixed-treatment occupation deletion. This is a robust aggregate direction for a specified tail contrast, not evidence of a monotone dose-response or a causal AI effect.

The agreement has sharp limits. One common-support interval includes zero under marginal inference, the predeclared simultaneous-band criterion does not establish all-six negativity, and mobility labels often disagree. Broad occupational assortativity explains nearly all of the realized excess mobility conflict over a stringent benchmark. The exploratory family decomposition places the negative conditional association more heavily on the Eloundou centroid, but its geometry is materially influenced by alpha and cannot rank families generally.

The durable conclusion is therefore neither that exposure choice is irrelevant nor that one index is uniquely correct. Alternative constructions can support the same aggregate employment tail contrast without being interchangeable for the rankings, comparisons, and economic interpretations attached to it. Robustness is statement-specific.

